\documentclass[../../main.tex]{subfiles}

\begin{document}

\section{Young, H\"older 与 Minkowski 不等式}\label{section:ch07-Young-Holder-Minkowski}

前一节对 $1\leqslant p\leqslant\infty$ 以及 $E$ 是实数的可测集, 我们引入线性空间 $L^{p}(E)$. 在 $p=1$ 与 $p=\infty$ 的情形, 我们定义了这些空间上的范数. 我们现在对 $1<p<\infty$ 定义 $L^{p}(E)$ 上的范数.

\begin{definition}[$L^{p}$ 范数]\label{definition:ch07-Lp范数}
对于 $E$ 为可测集, $1<p<\infty$, 以及 $L^{p}(E)$ 中的函数 $f$, 定义
\begin{align*}
\|f\|_{p}=\left[\int_{E}|f|^{p}\right]^{1/p}.
\end{align*}
\end{definition}

我们将证明泛函 $\|\cdot\|_{p}$ 是 $L^{p}(E)$ 上的一个范数. 事实上, 正齐次性显而易见. 此外, 根据第 4 章的命题 9, $\|f\|_{p}=0$ 当且仅当 $f$ 在 $E$ 上 a.e. 消失. 因此 $[f]$ 是线性空间 $L^{p}(E)\subseteq\mathcal{F}/\cong$ 的零元, 即 $f=0$. 剩下来要证明三角不等式, 即要证明对所有 $L^{p}(E)$ 中的 $f$, $g$
\begin{align*}
\|f+g\|_{p}\leqslant\|f\|_{p}+\|g\|_{p}.
\end{align*}
该不等式并不显而易见. 它称为 Minkowski 不等式.

\begin{definition}[共轭数]\label{definition:ch07-共轭数}
数 $p\in(1,\infty)$ 的共轭是数 $q=p/(p-1)$, 它是使得
\begin{align*}
\frac{1}{p}+\frac{1}{q}=1
\end{align*}
的唯一的数 $q\in(1,\infty)$. $1$ 的共轭定义为 $\infty$, 而 $\infty$ 的共轭定义为 $1$.
\end{definition}

\begin{proposition}[Young 不等式]\label{proposition:ch07-Young}
对于 $1<p<\infty$, $q$ 是 $p$ 的共轭, 以及任何两个正数 $a$ 和 $b$,
\begin{align*}
ab\leqslant\frac{a^{p}}{p}+\frac{b^{q}}{q}.
\end{align*}
\end{proposition}

\begin{proof}
函数 $e^{x}$ 有正的二阶导数, 因此是凸的, 即对任何数 $\lambda\in[0,1]$ 以及任何数 $u$ 和 $v$,
\begin{align*}
e^{\lambda u+(1-\lambda)v}\leqslant\lambda e^{u}+(1-\lambda)e^{v}.
\end{align*}
特别地, 设 $\lambda=1/p$, 因此 $1-\lambda=1/q$, 我们有
\begin{align*}
e^{\frac{1}{p}u+\frac{1}{q}v}\leqslant\frac{1}{p}e^{u}+\frac{1}{q}e^{v}.
\end{align*}
因此,
\begin{align*}
ab\leqslant\frac{a^{p}}{p}+\frac{b^{q}}{q}.
\end{align*}
\end{proof}

\begin{theorem}[H\"older 不等式]\label{theorem:ch07-Holder}
令 $E$ 为可测集, $1\leqslant p<\infty$, $q$ 为 $p$ 的共轭. 若 $f$ 属于 $L^{p}(E)$ 而 $g$ 属于 $L^{q}(E)$, 则它们的积 $fg$ 在 $E$ 上可积且 H\"older 不等式成立:
\begin{align}
\int_{E}|f\cdot g|\leqslant\|f\|_{p}\cdot\|g\|_{q}.\label{align:ch07-Holder}
\end{align}
此外, 若 $f\neq 0$, 函数\footnote{若 $f(x)\geqslant 0$ 函数 $\mathrm{sgn}(f)$ 取值 $1$, 而若 $f(x)<0$ 取值 $-1$. 由于 $f$ 在 $E$ 上 a.e. 有限, 因此在 $E$ 上 a.e. $\mathrm{sgn}(f)\cdot f=|f|$.} $f^{*}=\|f\|_{p}^{1-p}\cdot\mathrm{sgn}(f)\cdot|f|^{p-1}$ 属于 $L^{q}(E)$,
\begin{align}
\int_{E}f\cdot f^{*}=\|f\|_{p},\quad\text{且}\quad\|f^{*}\|_{q}=1.\label{align:ch07-共轭函数}
\end{align}
\end{theorem}

\begin{proof}
首先考虑 $p=1$ 的情形. 则 H\"older 不等式从积分的单调性和观察\eqref{align:ch07-Linfty本性上界}即 $\|g\|_{\infty}$ 是 $g$ 在 $E$ 上的本性上界得出. 观察到由于 $f^{*}=\mathrm{sgn}(f)$, \eqref{align:ch07-共轭函数}当 $p=1$, $q=\infty$ 时成立. 现在考虑 $p>1$. 假设 $f\neq 0$ 且 $g\neq 0$, 否则没什么要证的. 显然, 若 $f$ 用它的规范化 $f/\|f\|_{p}$ 代替且 $g$ 用它的规范化 $g/\|g\|_{q}$ 代替时 H\"older 不等式成立, 则它对 $f$ 和 $g$ 成立. 因此假设 $\|f\|_{p}=\|g\|_{q}=1$, 即
\begin{align*}
\int_{E}|f|^{p}=1\quad\text{以及}\quad\int_{E}|g|^{q}=1.
\end{align*}
在这种情形 H\"older 不等式成为
\begin{align*}
\int_{E}|f\cdot g|\leqslant 1.
\end{align*}
由于 $|f|^{p}$ 和 $|g|^{q}$ 在 $E$ 上可积, $f$ 和 $g$ 在 $E$ 上 a.e. 有限. 因此, 根据 Young 不等式,
\begin{align*}
\text{在 }E\text{ 上 a.e. }|f\cdot g|=|f|\cdot|g|\leqslant\frac{|f|^{p}}{p}+\frac{|g|^{q}}{q}.
\end{align*}
我们从积分的线性和比较判别法推出 $f\cdot g$ 在 $E$ 上可积, 且根据积分的单调性与线性, 有
\begin{align*}
\int_{E}|f\cdot g|\leqslant\frac{1}{p}\int_{E}|f|^{p}+\frac{1}{q}\int_{E}|g|^{q}=\frac{1}{p}+\frac{1}{q}=1.
\end{align*}
剩下来要证明\eqref{align:ch07-共轭函数}. 观察到
\begin{align*}
\text{在 }E\text{ 上 a.e. }f\cdot f^{*}=\|f\|_{p}^{1-p}\cdot|f|^{p}.
\end{align*}
因此
\begin{align*}
\int_{E}f\cdot f^{*}=\|f\|_{p}^{1-p}\cdot\int_{E}|f|^{p}=\|f\|_{p}^{1-p}\cdot\|f\|_{p}^{p}=\|f\|_{p}.
\end{align*}
由于 $q(p-1)=p$, $\|f^{*}\|_{q}=1$.
\end{proof}

对 $f\in L^{p}(E)$, $f\neq 0$, 称如上定义的函数 $f^{*}$ 为 $f$ 的共轭函数是很方便的.

\begin{theorem}[Minkowski 不等式]\label{theorem:ch07-Minkowski}
令 $E$ 为可测集而 $1\leqslant p\leqslant\infty$. 若函数 $f$ 和 $g$ 属于 $L^{p}(E)$, 则它们的和 $f+g$ 也如此, 此外,
\begin{align*}
\|f+g\|_{p}\leqslant\|f\|_{p}+\|g\|_{p}.
\end{align*}
\end{theorem}

\begin{proof}
在前面的内容中我们已考虑了 $p=1$ 和 $p=\infty$ 的情形, 所以这里考虑 $1<p<\infty$. 我们已从\eqref{align:ch07-三角不等式估计}推出了 $f+g$ 属于 $L^{p}(E)$. 假定 $f+g\neq 0$. 考虑 $(f+g)^{*}$, 即 $f+g$ 的共轭函数, 从积分的线性与 H\"older 不等式推出
\begin{align*}
\|f+g\|_{p}&=\int_{E}(f+g)\cdot(f+g)^{*}\\
&=\int_{E}f\cdot(f+g)^{*}+\int_{E}g\cdot(f+g)^{*}\\
&\leqslant\|f\|_{p}\cdot\|(f+g)^{*}\|_{q}+\|g\|_{p}\cdot\|(f+g)^{*}\|_{q}\\
&=\|f\|_{p}+\|g\|_{p}.
\end{align*}
\end{proof}

当 $p=q=2$ 时, H\"older 不等式的特殊情形有它自己的名字.

\begin{corollary}[Cauchy-Schwarz 不等式]\label{corollary:ch07-Cauchy-Schwarz}
令 $E$ 为可测集, 而 $f$ 和 $g$ 为 $E$ 上使得 $f^{2}$ 和 $g^{2}$ 在 $E$ 上可积的可测函数. 则它们的乘积 $fg$ 在 $E$ 上也是可积的, 且
\begin{align*}
\int_{E}|f\cdot g|\leqslant\sqrt{\int_{E}f^{2}}\cdot\sqrt{\int_{E}g^{2}}.
\end{align*}
\end{corollary}

\begin{corollary}\label{corollary:ch07-一致可积}
令 $E$ 为可测集而 $1<p<\infty$. 假定 $L^{p}(E)$ 中的函数簇 $\mathcal{F}$ 在以下意义下在 $L^{p}(E)$ 中有界: 存在常数 $M$, 使得
\begin{align*}
\text{对 }\mathcal{F}\text{ 中的所有 }f,\quad\|f\|_{p}\leqslant M.
\end{align*}
则簇 $\mathcal{F}$ 在 $E$ 上是一致可积的.
\end{corollary}

\begin{proof}
令 $\varepsilon>0$. 我们必须证明存在 $\delta>0$, 使得对 $\mathcal{F}$ 中的任何 $f$,
\begin{align*}
\text{若 }A\subseteq E\text{ 是可测的且 }m(A)<\delta,\quad\text{则}\int_{A}|f|<\varepsilon.
\end{align*}
令 $A$ 为 $E$ 的有限测度的可测子集. 考虑 $L^{p}(A)$ 和 $L^{q}(A)$, 其中 $q$ 是 $p$ 的共轭. 定义 $g$ 在 $A$ 上恒等于 $1$. 由于 $m(A)<\infty$, $g$ 属于 $L^{q}(A)$. H\"older 不等式应用于 $g$ 以及 $f$ 在 $A$ 上的限制推出,
\begin{align*}
\int_{A}|f|=\int_{A}|f|\cdot g\leqslant\left[\int_{A}|f|^{p}\right]^{1/p}\cdot\left[\int_{A}|g|^{q}\right]^{1/q}.
\end{align*}
但对所有 $\mathcal{F}$ 中的 $f$,
\begin{align*}
\left[\int_{A}|f|^{p}\right]^{1/p}\leqslant\left[\int_{E}|f|^{p}\right]^{1/p}\leqslant M,\quad\text{且}\left[\int_{A}|g|^{q}\right]^{1/q}=[m(A)]^{1/q}.
\end{align*}
因此, 对所有 $\mathcal{F}$ 中的 $f$,
\begin{align*}
\int_{A}|f|\leqslant M\cdot[m(A)]^{1/q}.
\end{align*}
于是, 对每个 $\varepsilon>0$, $\delta=[\varepsilon/M]^{q}$ 对应于 $\mathcal{F}$ 的一致可积准则的 $\varepsilon$ 挑战.
\end{proof}

\begin{corollary}\label{corollary:ch07-Lp包含}
令 $E$ 为有限测度的可测集, 而 $1\leqslant p_{1}<p_{2}\leqslant\infty$. 则 $L^{p_{2}}(E)\subseteq L^{p_{1}}(E)$. 此外,
\begin{align}
\text{对 }L^{p_{2}}(E)\text{ 中的所有 }f,\quad\|f\|_{p_{1}}\leqslant c\|f\|_{p_{2}},\label{align:ch07-Lp包含范数}
\end{align}
其中, 若 $p_{2}<\infty$ 有 $c=[m(E)]^{\frac{p_{2}-p_{1}}{p_{1}p_{2}}}$, 而若 $p_{2}=\infty$ 有 $c=[m(E)]^{\frac{1}{p_{1}}}$.
\end{corollary}

\begin{proof}
我们将 $p_{2}=\infty$ 的情形作为练习. 假设 $p_{2}<\infty$. 定义 $p=p_{2}/p_{1}>1$ 且令 $q$ 为 $p$ 的共轭. 令 $f$ 属于 $L^{p_{2}}(E)$. 观察到由于 $m(E)<\infty$, $f^{p_{1}}$ 属于 $L^{p}(E)$, 而 $g=\chi_{E}$ 属于 $L^{q}(E)$. 应用 H\"older 不等式, 则
\begin{align*}
\int_{E}|f|^{p_{1}}=\int_{E}|f|^{p_{1}}\cdot g\leqslant\|f^{p_{1}}\|_{p}\cdot\left[\int_{E}|g|^{q}\right]^{1/q}=\|f\|_{p_{2}}^{p_{1}}[m(E)]^{1/q}.
\end{align*}
每边取 $1/p_{1}$ 幂得到\eqref{align:ch07-Lp包含范数}.
\end{proof}

\begin{example}\label{example:ch07-Lp真子空间}
一般地, 对有限测度集 $E$ 和 $1\leqslant p_{1}<p_{2}\leqslant\infty$, $L^{p_{2}}(E)$ 是 $L^{p_{1}}(E)$ 的真子空间. 例如, 令 $E=(0,1]$, 而对 $0<x\leqslant 1$, $f$ 定义为 $f(x)=x^{\alpha}$, 其中 $-1/p_{1}<\alpha\leqslant-1/p_{2}$. 则 $f\in L^{p_{1}}(E)\sim L^{p_{2}}(E)$.
\end{example}

\begin{example}\label{example:ch07-Lp无包含}
一般地, 当 $E$ 是无限测度集时, $L^{p}(E)$ 空间之间没有包含关系. 例如, 对于 $E=(0,\infty)$, 而 $f$ 定义为
\begin{align*}
f(x)=\frac{x^{-1/2}}{1+|\ln x|},\quad x>0.
\end{align*}
$f$ 属于 $L^{p}(E)$ 当且仅当 $p=2$.
\end{example}

\begin{exercise}\label{exercise:ch07-06}
证明: 若 H\"older 不等式对规范化的函数成立, 则它对一般的函数也成立.
\end{exercise}

\begin{exercise}\label{exercise:ch07-07}
验证以上两个例子中关于 $L^{p}(E)$ 中函数 $f$ 的成员关系的断言.
\end{exercise}

\begin{exercise}\label{exercise:ch07-08}
令 $f$ 和 $g$ 属于 $L^{2}(E)$. 从积分的线性证明: 对任何数 $\lambda$,
\begin{align*}
\lambda^{2}\int_{E}f^{2}+2\lambda\int_{E}f\cdot g+\int_{E}g^{2}=\int_{E}(\lambda f+g)^{2}\geqslant 0.
\end{align*}
从这个平方和公式直接导出 Cauchy-Schwarz 不等式.
\end{exercise}

\begin{exercise}\label{exercise:ch07-09}
证明 Young 不等式成为等式当且仅当 $a^{p}=b^{q}=1$.
\end{exercise}

\begin{exercise}\label{exercise:ch07-10}
证明 H\"older 不等式成为等式当且仅当存在不全为零的常数 $\alpha$ 和 $\beta$, 使得在 $E$ 上 a.e.
\begin{align*}
\alpha|f|^{p}=\beta|g|^{q}.
\end{align*}
\end{exercise}

\begin{exercise}\label{exercise:ch07-11}
对 $\mathbf{R}^{n}$ 中的点 $x=(x_{1},x_{2},\cdots,x_{n})$, 定义 $T_{x}$ 为区间 $[1,n+1]$ 上的阶梯函数, 它在区间 $[k,k+1)$ 上取值 $x_{k}$, $1\leqslant k\leqslant n$. 对 $p\geqslant 1$, 定义 $\|X\|_{p}=\|T_{x}\|_{p}$, 即函数 $T_{x}$ 在 $L^{p}[1,n+1)$ 中的范数. 证明这定义了 $\mathbf{R}^{n}$ 上的范数. 对该范数叙述并证明 H\"older 与 Minkowski 不等式.
\end{exercise}

\begin{exercise}\label{exercise:ch07-12}
对 $1\leqslant p<\infty$ 和序列 $a=(a_{1},a_{2},\cdots)\in\ell^{p}$, 定义 $T_{a}$ 为区间 $[1,\infty)$ 上的函数, 它在 $[k,k+1)$ 上取值 $a_{k}$, $k=1,2,\cdots$. 证明 $T_{a}$ 属于 $L^{p}[1,\infty)$ 且 $\|a\|_{p}=\|T_{a}\|_{p}$. 以此叙述并证明 $\ell^{p}$ 中的 H\"older 与 Minkowski 不等式.
\end{exercise}

\begin{exercise}\label{exercise:ch07-13}
证明: 若 $f$ 是 $E$ 上属于 $L^{p_{1}}(E)$ 的有界函数, 则对任何 $p_{2}>p_{1}$, 它属于 $L^{p_{2}}(E)$.
\end{exercise}

\begin{exercise}\label{exercise:ch07-14}
证明: 若对 $x\in(0,1]$, $f(x)=\ln(1/x)$, 则对所有 $1\leqslant p<\infty$, $f$ 属于 $L^{p}(0,1]$ 但不属于 $L^{\infty}(0,1]$.
\end{exercise}

\begin{exercise}\label{exercise:ch07-15}
对三个函数的乘积叙述和证明 H\"older 不等式的延伸.
\end{exercise}

\begin{exercise}\label{exercise:ch07-16}
假定 $\{f_{n}\}$ 在 $L^{1}[0,1]$ 是有界的. $\{f_{n}\}$ 在 $[0,1]$ 上是一致可积的吗?
\end{exercise}

\begin{exercise}\label{exercise:ch07-17}
对于 $1\leqslant p<\infty$, 假定 $\{f_{n}\}$ 在 $L^{p}(\mathbf{R})$ 中是有界的. $\{f_{n}\}$ 是紧的吗?
\end{exercise}

\begin{exercise}\label{exercise:ch07-18}
假设 $m(E)<\infty$. 对 $f\in L^{\infty}(E)$, 证明 $\lim_{p\to\infty}\|f\|_{p}=\|f\|_{\infty}$.
\end{exercise}

\begin{exercise}\label{exercise:ch07-19}
对于 $1\leqslant p\leqslant\infty$, $q$ 是 $p$ 的共轭, 而 $f\in L^{p}(E)$, 证明:
\begin{align*}
\|f\|_{p}=\max_{g\in L^{q}(E),\|g\|_{q}\leqslant 1}\int_{E}f\cdot g.
\end{align*}
\end{exercise}

\begin{exercise}\label{exercise:ch07-20}
对于 $1\leqslant p\leqslant\infty$, $q$ 是 $p$ 的共轭, 而 $f\in L^{p}(E)$, 证明 $f=0$ 当且仅当
\begin{align*}
\text{对所有 }g\in L^{q}(E),\quad\int_{E}f\cdot g=0.
\end{align*}
\end{exercise}

\begin{exercise}\label{exercise:ch07-21}
对于 $1\leqslant p\leqslant\infty$, 找到参数 $\lambda$ 的值使得对所有 $f\in L^{p}[0,1]$,
\begin{align*}
\lim_{\alpha\to 0^{+}}\frac{1}{\varepsilon^{\lambda}}\int_{0}^{\varepsilon}f=0.
\end{align*}
\end{exercise}

\begin{exercise}\label{exercise:ch07-22}
(Riesz) 对于 $1\leqslant p<\infty$, 证明: 若 $[a,b]$ 上的绝对连续函数 $F$ 是一个 $L^{p}[a,b]$ 函数的不定积分, 则存在常数 $M>0$, 使得对 $[a,b]$ 的任何分划 $\{x_{0},\cdots,x_{n}\}$,
\begin{align*}
\sum_{k=1}^{n}\frac{|F(x_{k})-F(x_{k-1})|^{p}}{|x_{k}-x_{k-1}|^{p-1}}\leqslant M.
\end{align*}
\end{exercise}

\end{document}